% Esta seccion aloja tambien a secciones/objetivo_tecnicas.tex (issue #39),
% que continua con \subsection y se \input inmediatamente despues en main.tex.
% Se fusionaron para no tener dos secciones de "objetivo" con encabezado y
% parrafo de enlace propios.
%
% Fusiona lo que antes eran dos subsecciones (objetivo.tex e
% investigacion_obj.tex): ambas justificaban por que estudiar el ozono
% troposferico y se solapaban. Se conserva la definicion quimica, el hallazgo
% de la anomalia invernal en el AMM, los metodos de deteccion y el dato duro
% de la tendencia 1993-2014 del AMM frente a la CDMX; se recorto el parrafo
% generico de "crecimiento industrial y vehicular" que repetia, sin datos, lo
% que el parrafo de la tendencia 1993-2014 ya dice con cifras.

\section{Objetivo del análisis y técnicas estadísticas}
\label{sec:objetivo-tecnicas}

\subsection{Objeto de investigación}
\label{sec:objetivo}

El ozono \ce{O3} es una substancia elemental compuesta por tres átomos de oxígeno, y se encuentra tanto en la atmosfera superior (ozono estratosférico) como a nivel del suelo (ozono troposférico). El ozono estratosférico es producido de manera natural por la interacción entre el oxígeno molecular \ce{O2} y los rayos ultravioleta provenientes del sol, este uno de los principales responsables de atenuar la exposición de la superficie a los rayos ultravioleta, este efecto se considera sumamente importante, ya que su ausencia podría significar mayores riesgos para los seres vivos \parencite{Bernhard2023-lk}.

A diferencia del ozono estratosférico, el ozono troposférico es el principal ingrediente del \textit{smog} y se forma por transformaciones químicas entre los óxidos de nitrógeno \ce{NOx} y los compuestos orgánicos volátiles (COV) con el calor y la luz solar.

Áreas metropolitanas de Nuevo León con elevada actividad industrial han
presentado niveles de \ce{O3} y \ce{PM10} que superan los límites de la
NOM-020-SSA1-2021 \parencite{nom020ssa1} y la NOM-025-SSA1-2014
\parencite{nom025ssa1}. Estas tendencias se documentan desde hace más de
30 años, con distintas herramientas
de detección que permiten medir tanto las concentraciones de ozono (mediante
fotometría ultravioleta) como las de sus precursores, los óxidos de nitrógeno
\ce{NO} y \ce{NO2} (agrupados como \ce{NO_x}, medidos por quimioluminiscencia),
cuya relación fotoquímica es clave para entender la formación del contaminante
secundario \ce{O3} \parencite{SanfordSillman}.

Diversos estudios han evidenciado que el Área Metropolitana de Monterrey (AMM)
no es la excepción a este fenómeno: entre 1993 y 2014 se registró un aumento
sostenido de las concentraciones de ozono troposférico, en contraste con la
tendencia descendente observada en la Ciudad de México durante el mismo periodo,
lo cual ha sido atribuido al crecimiento industrial y vehicular de la región sin
un control equivalente sobre sus precursores atmosféricos, y a un papel
meteorológico adicionalmente modulado por el cambio climático
\parencite{atmos15070779}. Este comportamiento se agrava por las condiciones
topográficas del valle donde se asienta el AMM, rodeado de sierras que favorecen
el estancamiento de masas de aire durante episodios de alta presión atmosférica
(anticiclones), inversión térmica y baja velocidad de viento, factores que
reducen la dispersión de contaminantes y prolongan los periodos de exposición de
la población a concentraciones elevadas.
